% Main document additions: sections 4.3 -- 5.3
\documentclass[12pt]{article}
\usepackage[utf8]{inputenc}
\usepackage[T1]{fontenc}
\usepackage[vietnamese]{babel}
\usepackage{hyperref}
\usepackage{microtype}
\usepackage{enumitem}
\usepackage{parskip}
\begin{document}

% Force numbering to continue from section 4, subsection 2 -> next is 4.3
\setcounter{section}{4}
\setcounter{subsection}{2}

\subsection{4.3 Ghi nhận hành động do Automation thực hiện}
Trong quá trình phát triển, chúng tôi nhận thấy các hành động do engine tự động (automation engine) thực thi không được ghi lại vào hệ thống audit. Điều này gây ra hai vấn đề chính:
\begin{itemize}[noitemsep]
  \item Timeline trên Dashboard không phản ánh các thay đổi trạng thái do automation thực hiện.
  \item Thống kê (counts / analytics) tách biệt giữa hành động thủ công và hành động tự động, dẫn tới số liệu không chính xác.
\end{itemize}

Để khắc phục, backend (module \texttt{automation\_service.py}) được điều chỉnh để khi một \emph{rule} được kích hoạt và thực hiện hành động lên một thiết bị thì luôn tạo một bản ghi mới trong bảng \texttt{DeviceActivityLog} với các trường chính:
\begin{itemize}[noitemsep]
  \item \texttt{device\_id}, \texttt{action} (turn\_on / turn\_off / set\_level), \texttt{triggered\_by = 'automation\_rule'}
  \item \texttt{reason} chứa tên quy tắc hoặc diễn giải ngắn gọn về nguyên nhân
  \item \texttt{timestamp} là thời điểm thực thi
\end{itemize}

Hậu quả: Dashboard giờ đây hiển thị đầy đủ các mục lịch sử bao gồm cả hành động tự động, và các báo cáo thống kê (ví dụ số lần automation chạy trong khoảng thời gian) có thể sinh ra từ cùng một nguồn dữ liệu audit.

\subsection{4.4 Chỉnh sửa quy tắc Automation trên ứng dụng di động}
Đã mở rộng giao diện \texttt{AutomationRulesScreen.js} để hỗ trợ chế độ chỉnh sửa (edit mode). Điểm chính:
\begin{itemize}[noitemsep]
  \item Thêm state `editingRuleId` để phân biệt giữa tạo mới và chỉnh sửa.
  \item Form dùng chung được tiền điền (pre-fill) các giá trị hiện có của quy tắc khi vào chế độ edit: \texttt{rule\_name}, \texttt{conditions}, \texttt{action\_device\_id}, \texttt{action\_status}, \texttt{action\_level}.
  \item Gọi API \texttt{PUT /api/automation-rules/<rule\_id>} để cập nhật quy tắc thay vì luôn dùng POST.
\end{itemize}

Trong quá trình sửa, một lỗi ``Too many re-renders'' đã xuất hiện do một lời gọi cập nhật state vô tình nằm ngoài hàm reset; lỗi này được khắc phục bằng cách dọn dẹp các cập nhật không cần thiết và đảm bảo mọi \texttt{setState} chỉ xảy ra trong các handler sự kiện hợp lệ.

\subsection{4.5 Đồng bộ trạng thái thiết bị và chính xác khi toggling}
Vấn đề ban đầu: khi người dùng bật/tắt thiết bị trong Dashboard, giao diện hiển thị trạng thái tạm thời đúng nhưng sau khi rời màn hình, trạng thái quay lại giá trị trước đó. Nguyên nhân chính là luồng automation hoặc các tiến trình nền đã ghi đè trạng thái vừa thay đổi.

Các thay đổi thực hiện:
\begin{enumerate}[noitemsep]
  \item \textbf{Persist manual changes:} endpoint \texttt{POST /api/device-status} ghi trạng thái mới vào bảng \texttt{device} ngay lập tức.
  \item \textbf{Refresh sau khi toggle:} client (\texttt{DevicesScreen.js}) gọi reload danh sách thiết bị chính xác từ server sau khi nhận được response thành công, đảm bảo trạng thái UI đồng bộ với server.
  \item \textbf{Ghi log toggle:} chúng tôi bổ sung logging chi tiết cho response của thao tác toggle để dễ debug khi có xung đột với automation.
  \item \textbf{Automation tạo entry audit:} như mục 4.3, các hành động automation giờ sẽ xuất hiện trong \texttt{DeviceActivityLog} để có thể phân biệt hành vi tự động và thủ công.
\end{enumerate}

Kết quả: thao tác bật/tắt thủ công được duy trì và hiển thị nhất quán; nếu automation sau đó sửa đổi trạng thái thiết bị, hành động này sẽ được hiển thị như một sự kiện riêng biệt trong timeline.

\subsection{4.6 Xử lý timestamp và cải thiện phân tích Dashboard}
Phần lớn data activity ban đầu sử dụng timestamp ISO nhưng đôi khi là dạng 'naive' (không có ký hiệu timezone). Việc này dẫn tới lỗi hiển thị relative times ("7h ago" cho event vừa mới xảy ra) khi client giả định timezone khác.

Các biện pháp:
\begin{itemize}[noitemsep]
  \item Trên client, thêm helper \texttt{parseISOToDate()} để khi gặp timestamp không có hậu tố 'Z' sẽ mặc định coi đó là UTC, tránh sai lệch khi tính toán relative time.
  \item Dashboard aggregation (backend) nhận tham số \texttt{start\_date} từ client để sinh thống kê theo khoảng thời gian chính xác.
  \item Tách nguồn dữ liệu cho analytics: \texttt{automationActivities} và \texttt{onOffActivities} để dễ phân tích hành vi người dùng so với hành vi tự động.
\end{itemize}

% --- Section 5 ---
% Bắt đầu phần 5: đánh giá, thử nghiệm và triển khai
\setcounter{section}{5}
\setcounter{subsection}{0}

\subsection{5.1 Kiểm thử chức năng}
Đã tiến hành chuỗi kiểm thử bao gồm:
\begin{itemize}[noitemsep]
  \item Kiểm thử thủ công (manual QA): xác minh toggle thiết bị, chỉnh sửa quy tắc automation, xem timeline hiển thị automation run.
  \item Kiểm thử luồng end-to-end: từ mobile → API → automation engine → database → socket updates.
  \item Kiểm thử hồi quy (regression): đảm bảo những thay đổi ghi nhật ký không phá vỡ luồng toggle/thao tác người dùng.
\end{itemize}

Các trường hợp lỗi được phát hiện và sửa: render-loop trên màn hình AutomationRules, thiếu lưu \texttt{action\_device\_id} khi cập nhật quy tắc, và parsing timestamp gây sai relative time.

\subsection{5.2 Cập nhật tài liệu}
Đã cập nhật các tài liệu liên quan trong \texttt{docs/} để phản ánh:
\begin{itemize}[noitemsep]
  \item Hành vi mới của Dashboard (bao gồm automation activities trong timeline).
  \item API: chi tiết endpoint \texttt{PUT /api/automation-rules/{rule\_id}} và \texttt{POST /api/device-status}.
  \item Hướng dẫn vận hành: cách restart dịch vụ sau khi deploy, nơi xem logs và cách khôi phục database nếu cần.
\end{itemize}

\subsection{5.3 Các bước tiếp theo và đề xuất}
\begin{enumerate}[noitemsep]
  \item Thêm test tự động (unit/integration) cho \texttt{automation\_service} và API routes liên quan tới action logging.
  \item Thiết lập CI pipeline để tự động chạy test và (tùy chọn) build mobile app preview khi có PR.
  \item Xác minh môi trường sản xuất: cấu hình timezone thống nhất, job dọn logs định kỳ, và theo dõi alert khi automation chạy quá mức.
  \item Rà soát quyền (RBAC) cho các endpoint quản trị để tránh lạm dụng chức năng thay đổi trạng thái thiết bị.
\end{enumerate}

\end{document}
